% Source: https://artofproblemsolving.com/community/c3h1568113_mathcounts_questions
% Stable identity: aops-topic-1568113-q04
Since $n>0$ we have $x>y\ge1$. Write $n=(x-y)(x+y)=uv$ with $u=x-y\ge1$, $v=x+y$. Here $u$ and $v$ have the same parity, and $v-u=2y\ge2$. Conversely, any factorisation $n=uv$ with $u<v$ of the same parity gives $x=\dfrac{u+v}{2}$, $y=\dfrac{v-u}{2}$, and for $n<18$ these satisfy $x\le\dfrac{1+17}{2}=9<18$.

\textbf{Odd $n$.} Take $u=1$, $v=n$: this works for every odd $n\ge3$, giving $3,5,7,9,11,13,15,17$. ($n=1$ fails, since $1=1\cdot1$ has $u=v$.)

\textbf{$n\equiv2\pmod4$.} Then one of $u,v$ is even and the other odd, so these values ($2,6,10,14$) are impossible.

\textbf{Multiples of $4$.} Both factors must be even: $n=uv$ with $u=2u'$, $v=2v'$, $u'<v'$. This works for $8=2\cdot4$, $12=2\cdot6$, $16=2\cdot8$ but not for $4$ (only $2\cdot2$).

So the possible values are $3,5,7,8,9,11,12,13,15,16,17$, and their sum is
\[
(3+5+7+9+11+13+15+17)+(8+12+16)=80+36=\boxed{116}.
\]
